% ---------------------------------------------------------------------------
%  UNIT I -- Chapters 1 to 7
% ---------------------------------------------------------------------------
\chapter{Unit I --- Value Education, Self-exploration, Basic Human
  Aspirations and Harmony in the Self}
\chaptermark{Unit I --- Chapters 1 to 7}

\section{A. Understanding Value Education (Chapter 1)}

% ------------------------------------------------------------------------- 1
\Q{1}{What do you mean by values or human values? What is value education,
  and why is it needed?}{\tg{SEE}\mmk{6}}

\ans Character-oriented education that instils basic values and ethical value
in one's psyche is called \textbf{value-based education}. The subject that
enables us to understand `what is valuable' for human happiness is called
\textbf{value education}.

Values form the basis for all our thoughts, behaviour and actions. Once we
know what is valuable to us, these values become the basis, the anchor, for
our actions. Value education is important to help everyone improve the value
system he\,/\,she holds and put it to use. Once one has understood one's
values in life, one can examine and control the various choices one makes.

Value education enables us to understand our needs and visualise our goals
correctly, indicates the direction for their fulfilment, helps remove our
confusions and contradictions, brings harmony at all levels, and enables us to
rightly utilise technological innovations. We also need to understand the
universality of various human values, because only then can we have a definite
and common programme for value education --- and only then can we be assured
of a happy and harmonious human society.

% ------------------------------------------------------------------------- 2
\Q{2}{What is the need for value education? Explain its five
  aspects.}{\tg{SEE}\mmk{6}}

\ans
\begin{apts}
\item \textbf{Correct identification of our aspirations.} All human beings
      have aspirations and make plans. Before investing energy in actualising
      those plans, it is important to find out what we basically aspire for.
      Value education enables us to understand our needs, visualise our goals
      correctly and indicate the direction for their fulfilment.
\item \textbf{Understanding universal human values to fulfil our aspirations
      in continuity.} Identifying the aspiration is not enough; we must know
      how to fulfil it. Generally we pursue goals as per our appraisal and
      beliefs. Complete understanding of human values gives us a definite way
      to fulfil our aspirations.
\item \textbf{Complementarity of values and skills.} First we must know what
      really is conducive to human happiness --- happiness for one and for
      all, and at all times. This is the \emph{value domain}, the domain of
      wisdom. Secondly, we must learn methods and practices to actualise this
      goal in real life --- this is the \emph{domain of skills}. Both must go
      hand in hand.
\item \textbf{Evaluation of our beliefs.} In the absence of a correct
      understanding of universal human values we are driven by ad-hoc values
      and beliefs, which may or may not be true in reality. These beliefs come
      from what we read, see and hear, what our parents tell us, what
      magazines and TV say. They keep changing with time and are often
      mutually conflicting. Value education helps us evaluate our beliefs and
      assumed values.
\item \textbf{Technology and human values.} Technology is only a means to
      achieve what is considered valuable. It is not within the scope of
      technology to decide what is valuable --- that decision lies outside it.
      Without this decision technology can be aimless and can be put to any
      use, constructive or destructive. Hence technical education must be
      supplemented with value education.
\end{apts}

Value education is a crucial missing link in the present education system.
Because of this deficiency, most of our efforts may prove counterproductive,
and serious crises at the individual, societal and environmental level are
manifesting.

% ------------------------------------------------------------------------- 3
\Q{3}{What are the basic guidelines for value education?}{\tg{SEE}\mmk{6}}

\ans To qualify as an appropriate input in value education, the content of the
course must satisfy the following guidelines:

\begin{abul}
\item \textbf{Universal.} It must be applicable to all human beings
      irrespective of caste, creed, gender, nationality or religion, and be
      true at all times and all places.
\item \textbf{Rational.} It must appeal to human reasoning --- amenable to
      reasoning, and not based on dogmas or blind beliefs. It cannot be a set
      of sermons or do's and don'ts.
\item \textbf{Natural and verifiable.} It must be naturally acceptable to the
      human being who goes through the course, and living on the basis of such
      values must lead to fulfilment and happiness. It must be experientially
      verifiable and not based on dogmas, beliefs or assumptions.
\item \textbf{All-encompassing.} Value education is aimed at transforming our
      consciousness and living. Hence it must cover all the dimensions
      (thought, behaviour, work and realisation) and all the levels
      (individual, family, society, nature and existence) of human life and
      profession.
\item \textbf{Leading to harmony.} It must ultimately promote harmony within
      the individual, among human beings, and with the rest of nature.
\end{abul}

% ------------------------------------------------------------------------- 4
\Q{4}{What do you understand by the value of an entity? What is the value of
  a human being?}{\tg{SEE 2025}\mmk{4--8}}

\ans \textbf{The value of any unit in this existence is its participation in
the larger order of which it is a part.} For example, the value of a pen is
that it can write --- here writing is the participation of the pen in the
bigger order in which pen, paper and human being are all present. The value of
an eye is that it can be used for seeing. The value of a vegetable plant is
that it gives nutrition to animals and humans.

The \textbf{value of a human being} is the participation of the human being at
different levels in this order. This participation is seen in two forms ---
\emph{behaviour} and \emph{work}. The participation pertaining to behaviour
consists of the nine values in relationship: trust, respect, affection, care,
guidance, reverence, glory, gratitude and love. Working with material things,
we have two values: \textbf{utility value} and \textbf{artistic value}. All
these values are nothing but the participation of the human being in different
dimensions of living.

\begin{keybox}
\textbf{\sffamily Note for the 8-mark SEE version.} This question was set for
8 marks in June/July 2025 (Q2a). At that length, define the value of an
entity, give three examples, then give the value of a human being in
\emph{both} its forms --- the nine values in behaviour, named, and the two
values in work --- and close by stating that value is participation, not a
property the thing owns by itself.
\end{keybox}

% ------------------------------------------------------------------------- 5
\Q{5}{Illustrate the content of value education. What should its scope
  be?}{\mmk{5}}

\ans The scope of value education includes \textbf{all dimensions} (thought,
behaviour, work and realisation) and \textbf{all levels} (individual, family,
society, nature\,/\,existence) of human living.

Accordingly, the content of value education is: to understand myself, my
aspirations and my happiness; to understand the goal of human life
comprehensively; to understand the other entities in nature, the innate
interconnectedness and the co-existence in nature\,/\,existence; and finally
the role of the human being in this nature\,/\,existence entirely. Hence it
has to encompass the understanding of harmony at various levels ---
individual, family, society, nature and existence --- and finally, learning to
live in accordance with this understanding by being vigilant to one's thought,
behaviour and work.

% ------------------------------------------------------------------------- 6
\Q{6}{What do you mean by values? How do they differ from skills? How are
  values and skills complementary? Give two examples.}{\tg{SEE}\mmk{6}}

\ans \textbf{Values} mean importance or participation --- knowing what really
is conducive to human happiness, happiness for one and for all and at all
times. This is the \emph{value domain}, the domain of wisdom; it helps us
identify and set the right goals and proceed in the right direction.
\textbf{Skills} mean the qualities, training and capabilities --- the methods,
practices and techniques to actualise that goal in real life, in various
dimensions of human endeavour. This is the \emph{domain of skills}.

To fulfil our aspirations, both are necessary; hence there is an essential
complementarity between values and skills for the success of any human
endeavour.

\textbf{Example 1 (health).} I want to lead a healthy life. Only wishing for
good health will not keep my body fit; and without having understood the
meaning of health, I will not be able to choose things correctly to keep my
body fit. So I have to learn the skills to achieve the goal of good health ---
what food to consume, what physical workout to design. Without knowing the
meaning of good health (value), health cannot be achieved; and without the
skill, the goal cannot be reached.

\textbf{Example 2 (technology).} If we value the relationship with the
environment, we will work to create environment-friendly technologies (the
structure of technology) and also put them to right use (use of technology)
--- say for the enrichment of the environment and replenishment of natural
resources. The engineering skill produces the technology; the value decides
what to produce and how to use it. Conversely, if the relationship with the
environment is something we do not value, things could be the other way round.

% ------------------------------------------------------------------------- 7
\Q{7}{What is the need for value education in technical and other
  professional institutions?}{\mmk{5}}

\ans The present education system has become largely skill-based; the prime
emphasis is on science and technology, and on the skills to apply them. There
is very little in it that helps a student decide \emph{what is worth doing
with} those skills.

\begin{abul}
\item An engineer's decisions have consequences far beyond the engineer ---
      on other people, on society and on nature. Deciding what to produce,
      and how it should be used, is a question in the value domain, not the
      skill domain. Technology cannot answer it.
\item Without value education, professional competence is left aimless. The
      same competence can be put to constructive or destructive use, and
      nothing in the technical syllabus itself distinguishes the two.
\item Professionals with right understanding can ensure that their
      professional competence leads to mutual fulfilment in relationship and
      mutual prosperity with the rest of nature, rather than to exploitation
      of either.
\item The problems at the level of the individual, family, society and nature
      that we see today are a direct outcome of an incorrect understanding of
      happiness and prosperity, not of a shortage of technical skill. They
      cannot be solved by more skill alone.
\item Hence technical education must be \emph{supplemented} --- not replaced
      --- by value education, so that values and skills go hand in hand.
\end{abul}

\section{B. Self-exploration as the Process for Value Education (Chapter 2)}

% ------------------------------------------------------------------------- 8
\Q{8}{Define self-exploration. What is the content of
  self-exploration?}{\tg{CIE 1}\mmk{5}}

\ans \textbf{Self-exploration is the process of finding out what is valuable
to me by investigating within myself.} Since it is I who feels happy or
unhappy, successful or unsuccessful, whatever is right for me, true for me,
has to be judged within myself. Through self-exploration we get the value of
ourselves. We live with an entirety --- family, friends, air, soil, water,
trees --- and we want to understand our relationship with all these; for this
we need to start observing inside.

The main focus of self-exploration is \textbf{myself --- the human being}. Its
content is finding answers to the following two fundamental questions of all
human beings:

\begin{apts}
\item \textbf{Desire\,/\,Goal:} What is my (human) desire or goal? What do I
      really want in life, or what is the goal of human life?
\item \textbf{Programme:} What is my (human) programme for fulfilling the
      desire? What is the process to fulfil this basic aspiration?
\end{apts}

In short, these two questions cover the whole domain of human aspiration and
human endeavour; thus they form the content of self-exploration.

% ------------------------------------------------------------------------- 9
\Q{9}{Explain the process of self-exploration with a
  diagram.}{\tg{SEE}\mmk{6}}

\ans The process of self-exploration has to be kept in mind as follows:

\begin{abul}
\item \textbf{Whatever is being presented is a PROPOSAL.} Don't start by
      assuming it to be true or false.
\item \textbf{Verify it in your own right} --- not on the basis of scriptures,
      not on the basis of readings from instruments, and not on the basis of
      the assertion of other human beings.
\item \textbf{Firstly, verify the proposal on the basis of your natural
      acceptance.} Natural acceptance is a faculty present in each one of us;
      it is intact and invariant. We only have to start paying attention to
      it. (E.g. ask yourself: \emph{Is trust naturally acceptable to me in
      relationship, or is mistrust?} The answer comes from within,
      spontaneously.)
\item \textbf{Secondly, live according to the proposal to validate it
      experientially.} If the proposal is true, then in behaviour with other
      humans it will lead to \emph{mutual fulfilment\,/\,mutual happiness},
      and in work with the rest of nature it will lead to \emph{mutual
      prosperity}.
\end{abul}

\vspace{6pt}
\begin{center}
\begin{tikzpicture}[
  font=\sffamily\scriptsize,
  box/.style={draw=rulegrey,fill=panel,rounded corners=1pt,line width=0.7pt,
              align=left,inner sep=5pt,text width=36mm},
  head/.style={draw=none,fill=slate,text=white,rounded corners=1pt,
               align=center,inner sep=3.5pt,text width=36mm,font=\sffamily\bfseries\scriptsize},
  res/.style={draw=ocre,fill=tint,rounded corners=1pt,line width=0.8pt,
              align=center,inner sep=5pt,text width=124mm},
  ar/.style={-{Stealth[length=5pt]},line width=0.8pt,draw=ocre}]

\node[head] (h1) at (0,0) {PROPOSAL (input)};
\node[head] (h2) at (4.4,0) {VERIFY};
\node[head] (h3) at (8.8,0) {LIVE ACCORDINGLY};

\node[box,below=2pt of h1,text width=36mm] (b1)
  {It is a proposal\\[1pt]
   $\bullet$ Don't assume it true\\
   $\bullet$ Verify in your own right\\
   \hspace*{4pt}-- not from scriptures\\
   \hspace*{4pt}-- not from instruments\\
   \hspace*{4pt}-- not from others\\
   $\bullet$ Self-verification};
\node[box,below=2pt of h2,text width=36mm] (b2)
  {On the basis of\\[1pt]
   \textbf{Natural Acceptance}\\[2pt]
   Keep asking again and again:\\
   ``What is my natural acceptance?''};
\node[box,below=2pt of h3,text width=36mm] (b3)
  {\textbf{Experiential Validation}\\[2pt]
   $\bullet$ Behaviour with human\\
   \hspace*{4pt}$\rightarrow$ mutual happiness\\
   $\bullet$ Work with rest of nature\\
   \hspace*{4pt}$\rightarrow$ mutual prosperity};

\draw[ar] (b1.east) -- (b2.west);
\draw[ar] (b2.east) -- (b3.west);

\node[res] at (4.4,-3.55)
  {Results in \textbf{REALIZATION} and \textbf{UNDERSTANDING}\\[2pt]
   tested for \ \textbf{Assurance} \ $\vert$ \ \textbf{Satisfaction} \
   $\vert$ \ \textbf{Universality}\\
   (same for all Time, Space and Individual)};
\draw[ar] (b2.south) -- (4.4,-3.05);
\end{tikzpicture}
\end{center}

The answers we get on having realisation and understanding are:
\textbf{(a) Assuring} --- ``I am assured of the answer in myself'';
\textbf{(b) Satisfying} --- ``I am satisfied that the answers are fulfilling
for me''; \textbf{(c) Universal} --- the answers are the same for everyone,
invariant with respect to time, space and individual. If the answer fails any
of these criteria, it is most likely coming from past beliefs or conditioning
and not from natural acceptance, and needs to be re-verified.

% ------------------------------------------------------------------------ 10
\Q{10}{What do you mean by natural acceptance and experiential validation?
  Explain them as the two mechanisms of self-exploration.}{\tg{SEE}\mmk{5}}

\ans \textbf{Natural acceptance} implies unconditional and total acceptance of
the self, people and environment; it also refers to the absence of any
exception from others. It is a faculty present in each one of us, intact and
invariant, and we can always refer to it to know what is right for us. Once we
fully and truly commit ourselves on the basis of natural acceptance, we feel a
holistic sense of inner harmony, tranquillity and fulfilment.

\textbf{Experiential validation} is the process that infuses direct experience
with the learning environment and content. It may be regarded as a philosophy
and methodology in which the direct experience and focused reflection of the
individual helps to increase knowledge, develop skill and clarify values. We
live according to the proposal, and if it is true, our behaviour with humans
leads to mutual happiness and our work with the rest of nature leads to mutual
prosperity. When what we already believe to be true of us is validated by
situations, phenomena or outcomes, we may term it experiential validation.

Self-exploration takes place in \textbf{the Self, and not in the body}.

% ------------------------------------------------------------------------ 11
\Q{11}{Illustrate the purpose of self-exploration. (Seven points)}{\mmk{5}}

\ans
\begin{apts}
\item \textbf{It is a process of dialogue between ``what you are'' and ``what
      you really want to be''.} If these two are the same, there is no
      problem. If on investigation we find they are not the same, we are
      living with this contradiction, and we need to resolve this conflict
      within us.
\item \textbf{It is a process of self-evolution through self-investigation.}
      It successively enables us to bridge the gap between `what we are' and
      `what we want to be'; we become qualitatively better.
\item \textbf{It is a process of knowing oneself and, through that, knowing
      the entire existence.} The exploration starts by asking simple questions
      about ourselves, which gives clarity about our being, and then clarity
      about everything around us.
\item \textbf{It is a process of recognising one's relationship with every
      unit in existence and fulfilling it.} We become aware of our right
      relationship with other entities, discover the interconnectedness and
      co-existence, and live accordingly.
\item \textbf{It is a process of knowing human conduct, human character and
      living accordingly.} We all desire certainty and stability; once we know
      our own true nature, we understand our participation with the other
      things we live with --- this is the ethical or humane conduct.
\item \textbf{It is a process of being in harmony in oneself and in harmony
      with the entire existence.}
\item \textbf{It is a process of identifying our innateness and moving towards
      self-organisation and self-expression.}
\end{apts}

\begin{center}
{\sffamily\bfseries\color{slate}Svatva \ $\rightarrow$ \ Swatantrata \
 $\rightarrow$ \ Swarajya}
\end{center}

% ------------------------------------------------------------------------ 12
\Q{12}{What do you understand by the terms svatva, swatantrata and swarajya?
  How are they related?}{\mmk{5}}

\ans
\begin{abul}
\item \textbf{Svatva (innateness):} the innateness of the self --- the natural
      acceptance of harmony.
\item \textbf{Swatantrata (self-organisation):} being self-organised --- being
      in harmony with oneself.
\item \textbf{Swarajya (self-expression):} self-expression, self-extension ---
      living in harmony with others, and thus participation towards harmony in
      the whole existence.
\end{abul}

\begin{center}
{\sffamily\bfseries\color{slate}Svatva \ $\rightarrow$ \ Swatantrata \
 $\rightarrow$ \ Swarajya}
\end{center}

The \emph{svatva} is already there, intact in each one of us. By being in
dialogue with it, we attain \emph{swatantrata}, which enables us to work for
\emph{swarajya}. If we live in contradiction, it means we are not
self-organised and are living with pre-conditionings, having assumed certain
things and accumulated desires without first evaluating them --- then we are
\emph{partantra} (enslaved). On the other hand, when we identify our
innateness, what we really want to be, and establish a dialogue with it, it
enables us to start living with this harmony; it starts expressing itself
through our harmonious behaviour and work, and naturally extends to our
participation with the surroundings. This is working towards \emph{swarajya}.

% ------------------------------------------------------------------------ 13
\Q{13}{``Natural acceptance is innate, invariant and universal.'' Explain
  this statement with an example.}{\mmk{5}}

\ans We can verify proposals on the basis of the following characteristics of
natural acceptance:

\begin{apts}
\item \textbf{Natural acceptance does not change with time.} It remains
      invariant with time. For example, our natural acceptance for trust and
      respect does not change with age; people a hundred years ago also had
      the same natural acceptance.
\item \textbf{It does not depend on the place.} Whether we are in New Delhi,
      New York or Abu Dhabi, if we address our natural acceptance the answer
      would still be the same.
\item \textbf{It does not depend on our beliefs or past conditionings.} No
      matter how deep our belief or past conditioning, as long as we ask
      ourselves the question sincerely and refer deep within ourselves, the
      answer will always be the same.
\item \textbf{This natural acceptance is `constantly there', something we can
      refer to.} Try it: think of cheating or exploiting someone. The moment
      you think of this, you sense a contradiction within and feel unhappy
      that very instant. Every time we do something not acceptable to us,
      there is a contradiction in us, because the thought or deed conflicts
      with our own natural acceptance.
\item \textbf{Natural acceptance is the same for all of us.} It is part and
      parcel of every human being; it is part of human-ness. No human being
      finds disrespect acceptable in relationship. Our assumptions and
      choices, likes and dislikes may be different, but on basic and common
      issues like the need to be happy, to be respected, to be prosperous, we
      are all the same.
\end{apts}

\textbf{Example.} Take the proposal --- `respect is a value in human
relation'. When I verify at the level of natural acceptance, I find that it is
naturally acceptable to me. Similarly, when I behave with respect, it is
mutually fulfilling to me and to the other. Thus the proposal is true. If it
fails either of the two tests, it is untrue.

\section{C. The Basic Human Aspirations (Chapter 3)}

% ------------------------------------------------------------------------ 14
\Q{14}{What is happiness? ``To be in a state of harmony is happiness'' ---
  explain and illustrate with two examples from your day-to-day
  life.}{\tg{SEE}\mmk{6}}

\ans Happiness may be defined as \textbf{being in harmony\,/\,synergy in the
state or situation that I live in}.

\begin{keybox}
``The state\,/\,situation in which I live, if there is harmony\,/\,synergy in
it, then I like to be in that state\,/\,situation.'' \quad i.e. \quad
\textbf{To be in a state of liking is happiness.}

\vspace{3pt}
Conversely: ``The state\,/\,situation in which I live, if there is
conflict\,/\,contradiction in it, then I do not like to be in that state or
situation'' --- \textbf{to be in a state of disliking is unhappiness}.
\end{keybox}

Thus, \textbf{to be in a state of harmony is happiness; to be in a state of
disharmony or contradiction is unhappiness.} When we are in such a state of
happiness we experience no struggle, no contradiction or conflict within, and
we wish to have its continuity. One important characteristic of this feeling
is that \emph{we like to continue it}.

\textbf{Example 1 --- Respect.} Respect is a state of harmony between two
human beings. When I respect the other and the other respects me, I like to be
in that situation; it gives me happiness.

\textbf{Example 2 --- Harmony in my own thoughts.} When there is harmony in my
thoughts and feelings, I feel relaxed and want to be in that situation. This
feeling is happiness. If this harmony is disturbed --- say I am confused or in
conflict about a decision --- I feel uneasy, i.e. unhappy.

It should be noted that we do get an impression of happiness through sensory
interaction (tasty food, a beautiful picture, a sweet fragrance). However,
these impressions are always short-lived and their continuity can never be
ensured.

% ------------------------------------------------------------------------ 15
\Q{15}{What is the meaning of prosperity? How can you say that you are
  prosperous?}{\mmk{5}}

\ans \textbf{Prosperity is the feeling of having or making available more than
the required physical facilities.} Physical facilities are the material things
we use to fulfil the needs of the body; when we are able to cater to the needs
of the body adequately, we feel prosperous.

To ascertain prosperity, \textbf{two things} are essential:

\begin{apts}
\item \textbf{Correct assessment of the need for physical facilities} ---
      identification of the required quantity. We can be prosperous only if
      there is a limit to the need for physical facilities. If there is no
      limit, then whatsoever be the availability with us, the feeling of
      prosperity cannot be assured.
\item \textbf{The competence of making available more than the required
      physical facilities} (through production). Just assessing the need is
      not enough; we need to be able to produce or make available more than
      the perceived need.
\end{apts}

Since for all our physical facilities we are directly or indirectly dependent
on nature, the continuity of prosperity can be ensured only if our production
systems are in harmony with nature.

% ------------------------------------------------------------------------ 16
\Q{16}{How do prosperity and wealth differ? What is more acceptable to us and
  why?}{\tg{CIE 1}\tg{SEE}\mmk{8}}

\ans
\begin{cmptable}{Wealth}{Prosperity}
Wealth is a \textbf{physical thing}. It means having money, or having a lot of
  physical facilities, or both. &
Prosperity is a \textbf{feeling} of having more than required physical
  facilities. It is not just physical facilities. \\ \hline
It can be quantified and accumulated without limit; there is no point at which
  it is ``enough''. &
It requires (a) correct assessment of the need, and (b) the competence to make
  available more than that need. \\ \hline
Having wealth does not by itself lead to sharing; a person may have a lot of
  money but not want to share even a bit of it. &
The feeling of prosperity naturally leads to sharing with the other, becoming
  an aid by enriching the other. Deprivation leads to exploiting the other ---
  this is a simple test of prosperity. \\ \hline
Wealth is only a \textbf{part} of prosperity --- a means to the end. &
Prosperity is the \textbf{end}; the basic desire is to feel prosperous, to
  have a feeling of having enough. \\ \hline
\end{cmptable}

\textbf{Illustration.} A person has a lot of wealth, but gets worried when a
guest may stay longer than expected and might have to be fed. The person
\emph{has} wealth, but \emph{feels} deprived. On the other hand, someone who
does not have a lot of wealth may welcome you into his\,/\,her house and ask
you to stay back for a few days and help out --- this is an indication of
feeling prosperous. If one felt prosperous, one would have shared what one
has, since there is far more than enough anyway.

\textbf{What is more acceptable and why.} On checking our natural acceptance,
we find that \textbf{prosperity is more acceptable}. Given the choice between
\emph{accumulating more and more wealth while feeling deprived} and
\emph{having requisite wealth and feeling prosperous}, we find the latter
naturally acceptable. Not only do we want wealth, we want to feel prosperous
too. Wealth is only a means to that end. In order to feel prosperous, we need
to first decide how much wealth is needed, else it is like trying to fill
water in a glass that has no bottom --- the glass will never be filled,
howsoever one may try. This is why today we are busy accumulating wealth but
do not feel prosperous: we do not identify our need, and hence no matter how
much we have, it is always less, and we feel deprived.

% ------------------------------------------------------------------------ 17
\Q{17}{Critically examine the prevailing notions of happiness and prosperity
  and their consequences. \emph{(Also set as:} There are many problems
  manifest today at the level of individual, family, society and nature.
  Identify some of these.\emph{)}}{\tg{CIE 2 QB}\mmk{5}}

\ans In the current scenario, we are generally trying to achieve happiness and
prosperity by maximising \textbf{accumulation and consumption of physical
facilities}. This is an attempt to achieve happiness through pleasant sensory
interactions. Physical facilities are not seen in terms of fulfilling bodily
needs but as a means of maximising happiness. This has resulted in a wrong
assessment of wants for physical facilities as being unlimited.

But this pursuit is self-defeating. Neither can we hope to achieve continuous
happiness through sensory interactions, nor can we have prosperity, as it
amounts to trying to fulfil unlimited wants through limited resources. This
effort is engendering problems at all levels; it is becoming anti-ecological
and anti-people, and is threatening human survival itself. Some consequences:

\begin{abul}
\item \textbf{At the level of the individual} --- rising problems of
      depression, psychological disorders, suicides, stress, insecurity,
      psycho-somatic diseases, loneliness, lack of confidence and conviction.
\item \textbf{At the level of the family} --- breaking of joint families,
      mistrust, conflict between older and younger generations, insecurity in
      relationships, divorce, dowry tortures, family feuds, wasteful
      expenditure in family functions, neglect of older people.
\item \textbf{At the level of society} --- growing incidences of terrorism and
      naxalism, rising communalism, spreading casteism, racial and ethnic
      struggle, wars between nations, attempts at genocide, fear of nuclear
      and genetic warfare, corruption, adulteration, proliferation of lethal
      weapons.
\item \textbf{At the level of nature} --- global warming; water, air, soil and
      noise pollution; resource depletion of minerals and mineral oils;
      sizeable deforestation; loss of fertility of soil.
\end{abul}

All these problems are a direct outcome of an \textbf{incorrect
understanding} --- our wrong notion about happiness and prosperity and their
continuity. It therefore calls for an urgent need for human beings to
correctly understand happiness and prosperity as well as the sustainable way
to achieve these.

% ------------------------------------------------------------------------ 18
\Q{18}{What do the abbreviations SVDD, SSDD and SSSS signify?}{\mmk{4}}

\ans To achieve our basic aspirations we need to work for right understanding
as the base, on which we can work for relationship and then physical
facilities. Today we are not working according to this, and hence we find two
kinds of people in the world:

\begin{apts}
\item \textbf{SVDD --- Sadhan Viheen Dukhi Daridra:} Materially deficient,
      unhappy and deprived. Those who do not have physical
      facilities\,/\,wealth and feel unhappy and deprived.
\item \textbf{SSDD --- Sadhan Sampann Dukhi Daridra:} Materially affluent,
      unhappy and deprived. Those who have physical facilities\,/\,wealth but
      still feel unhappy and deprived.
\end{apts}

These are states we do not want to be in. We want to move to the third
category:

\begin{keybox}
\textbf{SSSS --- Sadhan Sampann Sukhi Samriddha:} Materially adequate, happy
and prosperous.
\end{keybox}

Presently, as we look around, we find most people in the first two categories,
while the natural acceptance of all human beings is to be in the category of
\textbf{SSSS}.

\section{D. The Programme to Fulfil Basic Human Aspirations (Chapter 4)}

% ------------------------------------------------------------------------ 19
\Q{19}{What are the requirements to fulfil basic human aspirations? Explain
  each and give the correct priority among them.}{\tg{CIE 2 QB}\mmk{5}}

\ans Our basic aspirations are \textbf{happiness (mutual fulfilment)} and
\textbf{prosperity (mutual prosperity)}. Happiness is ensured by relationships
with other human beings, and prosperity is ensured by working on physical
facilities. Three things are needed:

\begin{apts}
\item \textbf{Right Understanding.} This refers to the higher-order human
      skills --- the need to learn and utilise our intelligence most
      effectively. It is needed in myself; it is a need of `I'.
\item \textbf{Relationship.} This refers to the interpersonal relationships a
      person builds in his\,/\,her life --- at home, at the workplace and in
      society. We want to have the right feelings in relationship; this too is
      a need of `I'.
\item \textbf{Physical Facilities.} This includes the physiological needs of
      individuals and indicates the necessities as well as the comforts of
      life. It means the feeling of having or being able to have more physical
      facilities than is needed. This is a need of the Body.
\end{apts}

\begin{keybox}
\textbf{Correct priority: \ Right Understanding \ $\rightarrow$ \
Relationship \ $\rightarrow$ \ Physical Facilities.}
\end{keybox}

Right understanding comes first, because in order to resolve the issues in
human relationships we need to understand them first, and this comes from
`right understanding of relationship'. Similarly, in order to be prosperous
and to enrich nature, we need the right understanding --- it enables us to
work out our requirements for physical facilities and hence correctly
distinguish between wealth and prosperity. With nature as well, we need to
understand the harmony in nature and how we can complement it.

% ------------------------------------------------------------------------ 20
\Q{20}{``Right understanding $+$ Relationship $=$ Mutual fulfilment; Right
  understanding $+$ Physical facilities $=$ Mutual prosperity.'' Illustrate
  the above with two examples for each.}{\tg{CIE 1}\mmk{8}}

\ans Our basic aspirations are happiness (mutual fulfilment) and prosperity
(mutual prosperity). Happiness is ensured by relationships with other human
beings; prosperity is ensured by working on physical facilities. The reason we
are unable to have either today is that we have not focused on one more
aspect --- \textbf{right understanding}.

\begin{cmptable}{Right understanding $+$ Relationship $=$ Mutual fulfilment}%
                {Right understanding $+$ Physical facilities $=$ Mutual
                 prosperity}
With right understanding, we can be happy in ourselves and work to have
  fulfilling relationships with humans --- both sides feel fulfilled and
  satisfied from the interaction. &
With right understanding we correctly identify our need for physical
  facilities, produce more than required, and interact with nature in a way
  that is mutually enriching. \\ \hline
\end{cmptable}

\textbf{Examples for Mutual fulfilment:}
\begin{apts}
\item \textbf{Trust in the family.} With right understanding, I see that at
      the level of intention every human being wants my happiness, and only
      competence may be lacking. So when a family member fails to do something
      for me, I do not doubt his\,/\,her intention; I become a help. The
      relationship is fulfilling to both --- mutual fulfilment. Without right
      understanding, I doubt the intention, we fight, and both of us are
      unhappy.
\item \textbf{Respect for a classmate.} With right understanding, I evaluate
      the other rightly --- at the level of `I' the other is similar to me in
      aspiration, programme and potential. I neither over-evaluate nor
      under-evaluate. The other feels rightly evaluated and I feel at ease;
      both are fulfilled. Without it, I differentiate on body, wealth or
      beliefs, and both feel uneasy.
\end{apts}

\textbf{Examples for Mutual prosperity:}
\begin{apts}
\item \textbf{Farming.} With right understanding, I use a farming method that
      retains\,/\,enhances the fertility of the soil. I keep getting grain to
      feed my body and the soil is enriched --- both the human and nature
      gain. Without it, heavy chemical use makes the land unfit for
      agriculture; I lose and so does nature.
\item \textbf{Wood and trees.} The wood one person requires in a lifetime can
      be obtained from about four full-grown trees; a person can plant far
      more than four in a lifetime --- ten, twenty or a hundred. With right
      understanding I identify my limited need and plant more than I take, so
      I feel prosperous and nature is enriched. Without it, I over-produce and
      deforest, feeling deprived all the same.
\end{apts}

\begin{keybox}
Right Understanding $\rightarrow$ Relationship $\rightarrow$ \textbf{Happiness
with human beings} \hfill \textbar\hfill Right Understanding $\rightarrow$
Physical Facilities $\rightarrow$ \textbf{Mutual prosperity with the rest of
nature}
\end{keybox}

% ------------------------------------------------------------------------ 21
\Q{21}{``Physical facilities are necessary and complete for animals, while
  they are necessary but not complete for humans.'' Comment. \emph{(Also set
  as:} For human being physical facility is necessary, but relationship is
  also necessary --- justify with an example comparing with animal
  order.\emph{)}}{\tg{SEE 2025}\mmk{5--6}}

\ans
\begin{cmptable}{For Animals --- necessary and complete}%
                {For Humans --- necessary but not complete}
Animals need physical things to survive, mainly to take care of their body. A
  cow will look for food when it is hungry; once it gets the grass or fodder,
  it eats, and then sits around to chew at leisure. As long as animals have
  physical things they are largely fine. They don't desire other things like
  knowledge, or a peaceful animal society, or getting a good MBA. &
While physical facilities are necessary for human beings, they are not
  complete by themselves to fulfil our needs. Once you have had your fill, you
  do not just sit around and relax. We all have other needs and other plans
  --- going to a movie, reading a book, going to college, spending time with
  family and friends --- the list is endless. \\ \hline
\end{cmptable}

The reason is that a human being is the \textbf{co-existence of the Self (`I')
and the Body}. Physical facilities are needed for the Body; having them
ensures the fulfilment of the need of the Body. It does not address the need
of `I' --- happiness, trust, respect. Hence, besides physical facilities, we
also want \textbf{relationship} (mutual fulfilment) and \textbf{right
understanding}. Our needs are more than just physical facilities; we need
physical facilities, but the need does not end there.

% ------------------------------------------------------------------------ 22
\Q{22}{What do you mean by animal consciousness and human consciousness?
  Explain with the help of a diagram. How does the transformation take
  place?}{\mmk{5}}

\ans
\begin{abul}
\item \textbf{Animal Consciousness:} giving all priority to physical
      facilities only, or to live solely on the basis of physical facilities.
\item \textbf{Human Consciousness:} living with all three --- Right
      understanding, Relationship and Physical facilities.
\end{abul}

\vspace{4pt}
\begin{center}
\begin{tikzpicture}[
  font=\sffamily\scriptsize,
  bx/.style={draw=rulegrey,fill=panel,rounded corners=1pt,line width=0.7pt,
             align=center,inner sep=5pt},
  hi/.style={draw=ocre,fill=tint,rounded corners=1pt,line width=0.9pt,
             align=center,inner sep=5pt},
  ar/.style={-{Stealth[length=5pt]},line width=0.8pt,draw=ocre}]

\node[bx] (ru) at (0,0) {\textbf{Right}\\\textbf{Understanding}};
\node[bx] (rel) at (4.3,0.85) {Relationship with human\\$\rightarrow$
  \textbf{Mutual Happiness}};
\node[bx] (pf) at (4.3,-0.85) {Physical facilities\\$\rightarrow$
  \textbf{Mutual Prosperity}};
\draw[ar] (ru.east) -- (rel.west);
\draw[ar] (ru.east) -- (pf.west);

\node[bx,text width=52mm] (ac) at (0.4,-2.65)
  {To live \emph{solely} on the basis of physical facilities\\
   $=$ \textbf{ANIMAL CONSCIOUSNESS}};
\node[hi,text width=52mm] (hc) at (0.4,-4.55)
  {To live with \emph{all three}\\
   $=$ \textbf{HUMAN CONSCIOUSNESS}\\
   {\scriptsize this is where we want to be}};
\draw[ar] (ac.south) -- node[right=2pt,font=\sffamily\scriptsize\bfseries,
  text=ocre]{transformation} (hc.north);
\end{tikzpicture}
\end{center}

Physical facilities: for humans --- \textbf{necessary but not complete}; for
animals --- \textbf{necessary and complete}.

From the diagram we can say: for animals, physical facilities are necessary as
well as complete, whereas for human beings they are necessary but not
complete. Working only for physical facilities is living with Animal
Consciousness. Working for right understanding as the first priority, followed
by relationship and physical facilities, implies living with Human
Consciousness.

\textbf{The transformation} from Animal Consciousness to Human Consciousness
can be accomplished only by working for \textbf{right understanding as the
first priority}. This transformation forms the basis for human values and
value-based living. The content of education is the understanding of harmony
at all four levels of our existence --- from myself to the entire existence;
and \emph{sanskar} (right living) refers to the ability to live in harmony at
all four levels. This dimension of society works to ensure `right
understanding' and `right feelings' in every individual, i.e. the
all-encompassing solution called \emph{samadhan}, and ensures that our
succeeding generations have both the content and the environment available to
work towards achieving their goal of continuous happiness and prosperity.

\section{E. Human Being as Co-existence of Self (`I') and Body (Chapter 5)}

% ------------------------------------------------------------------------ 23
\Q{23}{Distinguish between the needs of the Self and the needs of the
  Body.}{\tg{SEE}\mmk{5}}

\ans A human being is the \textbf{co-existence of the Self (`I') and the
Body}, with an exchange of information between the two. Their needs are
fundamentally different:

\renewcommand{\arraystretch}{1.3}
\arrayrulecolor{rulegrey}
\begin{longtable}{P{22mm} P{63mm} P{63mm}}
\rowcolor{slate}\hdr{} & \hdr{Self (`I')} & \hdr{Body} \\ \hline \endfirsthead
\rowcolor{slate}\hdr{} & \hdr{Self (`I')} & \hdr{Body} \\ \hline \endhead
\textbf{Needs are} & Trust, respect, happiness\ldots\ i.e.
  \textbf{Happiness (sukha)} &
  Food, clothing, shelter, instruments\ldots\ i.e. \textbf{Physical Facilities
  (suvidha)} \\ \hline
\textbf{In time} & \textbf{Continuous} --- I want to be happy and respected
  all the time, not even for a moment otherwise. &
  \textbf{Temporary} --- food, clothes, shelter are needed only periodically.
  Even breathing is not continuous; we inhale, then exhale. \\ \hline
\textbf{In quantity} & \textbf{Qualitative} (no quantity) --- we cannot speak
  of one kg of respect or two litres of affection. Either the feeling is there
  or it is not. &
  \textbf{Quantitative} (limited in quantity) --- four chapattis, one bicycle.
  Beyond the limit, consumption becomes \emph{necessary and tasteful}
  $\rightarrow$ \emph{unnecessary but tasty} $\rightarrow$ \emph{unnecessary
  and tasteless} $\rightarrow$ \emph{intolerable}. \\ \hline
\textbf{Fulfilled by} & \textbf{Right understanding and right feelings} &
  \textbf{Physico-chemical things} --- food, clothing, etc. \\ \hline
\textbf{Activities} & Desiring, thinking, etc. --- knowing, assuming,
  recognising, fulfilling &
  Breathing, heartbeat, etc. --- recognising, fulfilling \\ \hline
\textbf{Type} & \textbf{Conscious} (non-material) &
  \textbf{Physico-chemical} (material) \\ \hline
\end{longtable}

\textbf{The confusion today.} A common mistake is that we mix these two sets
of needs and assume that ``all we need is physical facilities (\emph{suvidha}),
and that will automatically ensure happiness (\emph{sukha})''. The reality is
that we need both, since one is the need of the Body and the other the need of
`I'. \textbf{One cannot replace the other}, and the programmes for the two are
also different. For example, sitting in an air-conditioned room on a
comfortable sofa with a person for whom you have a feeling of opposition ---
the body is well taken care of, but you are unhappy. Our gross
misunderstandings today are: \textbf{Body $=$ `I'; \ Facilities $=$ Happiness;
\ Clothes $=$ Respect.}

\section{F. Harmony in the Self --- and with the Body (Chapters 6 \& 7)}

% ------------------------------------------------------------------------ 24
\Q{24}{How can we ensure harmony in the Self (`I')?}{\tg{CIE 1}\mmk{4}}

\ans The activities of \textbf{imaging (desire)}, \textbf{analysing (thought)}
and \textbf{selecting\,/\,tasting (expectation)} are constantly taking place
in `I'; together they are called \textbf{imagination}. Today these are largely
set either by \textbf{pre-conditioning (manyata)} or by \textbf{sensations
from the body} --- i.e. from the outside --- and are not self-verified on the
basis of our natural acceptance. In this state we are \emph{partantra}
(enslaved), and our desires, thoughts and expectations are in conflict among
themselves and with our own natural acceptance. This is the cause of
unhappiness.

The way to ensure harmony in the Self is a \textbf{four-step process}:

\begin{apts}
\item Becoming aware that the human being is the \textbf{co-existence of `I'
      and the Body}.
\item Becoming aware that the \textbf{Body is only an instrument of `I'}; `I'
      is the seer, doer and enjoyer.
\item Becoming aware of the activities of \textbf{desire, thought and
      expectation}, and passing each of these through your \textbf{natural
      acceptance}.
\item Understanding the \textbf{harmony at all levels of our existence} by
      verifying the proposals at the level of natural acceptance. This leads
      to \textbf{realization and understanding}, which in turn become the
      basis for desire, thought and expectation --- and this leads to harmony
      in `I' in continuity.
\end{apts}

\textbf{The outcome:} desires, thoughts and expectations become definite and
have a clear flow, so there is no contradiction; we have clarity about
ourselves, our basic aspiration and the way to fulfil it; we understand all
the levels of our living and live accordingly; and we live in a state of
\emph{svatantrata} --- self-organised in our imagination, behaviour and work.
This results in continuous happiness and prosperity.

\begin{keybox}
Realization \& Understanding (1, 2) $\rightarrow$ Desire (3) $\rightarrow$
Thought (4) $\rightarrow$ Expectation (5) \ $=$ \ \emph{svatantra}
(self-organised)
\end{keybox}

% ------------------------------------------------------------------------ 25
\Q{25}{Define Sanyama and Svasthya. How are the two related?}{\mmk{4--5}}

\ans The harmony of `I' with the Body is in the form of \textbf{Sanyama} on
the part of `I' and \textbf{Svasthya} on the part of the Body.

\begin{abul}
\item \textbf{Sanyama (Self-regulation):} the feeling of responsibility in the
      Self (`I') for \textbf{nurturing, protecting and rightly utilising} the
      Body.
\item \textbf{Svasthya (Health):} the state of the Body when it is \textbf{fit
      to act according to the needs of the Self (`I')}, and there is
      \textbf{harmony among the different parts of the Body}.
\end{abul}

\textbf{Relation: Sanyama is the basis of Svasthya.} When I live with Sanyama,
there is harmony among the different parts of the Body and the Body is fit to
be used as an instrument of `I'. If there is Sanyama, health can be ensured;
if Sanyama is not there, even good health can be lost. Today, in place of
being responsible to the Body, we rely more on medication; we go for newer
sources of sensual pleasure and irresponsible practices in living, and thus
produce new diseases day by day. An increase in hospitals or medical grants is
no substitute for Sanyama.

% ------------------------------------------------------------------------ 26
\Q{26}{``Realization and understanding are essential for happiness and
  harmony.'' Explain.}{\tg{SEE}\mmk{5}}

\ans Happiness is being in a state of harmony; unhappiness is a state of
conflict and contradiction. We have seen that this conflict is primarily
\textbf{inside us} --- between our imagination (desire, thought, expectation)
and our natural acceptance.

As long as our desires are being set from the outside --- by a sensation from
the body or by a pre-conditioning (\emph{manyata}) --- there is every chance
that we are in conflict, and we are \emph{partantra} (enslaved). Not only are
the desires, thoughts and expectations in conflict among themselves, they are
also in conflict with our own natural acceptance, and this creates
unhappiness.

Through the process of self-exploration, the activities of \textbf{realization
and understanding} get activated. Once we start operating at the level of
realization and understanding, our desires, thoughts and expectations get
aligned with our own natural acceptance, and we become \emph{svatantra}
(self-organised). There is self-organisation in my activities, leading to
continuity of happiness --- this is \textbf{harmony in the Self (`I')}.

Further, since the understanding is invariant, the desires arising from it are
also definite, and hence the thoughts, selections, behaviour and work are also
in harmony. This is what gives \textbf{definiteness of human conduct}.
Realization and understanding are therefore not optional extras --- they are
the very basis of continuous happiness, which is the basic human aspiration.
